% 消融实验结果表格 - LaTeX格式
% 使用方法: 复制到论文的table环境中

% ===== 表格1: 完整性能对比 =====
\begin{table}[htbp]
\centering
\caption{Ablation Study Results on Town01 Dataset (5000 frames)}
\label{tab:ablation_full}
\begin{tabular}{l|ccc|cc}
\hline
\textbf{Configuration} & \textbf{VT} & \textbf{Exp.} & \textbf{RMSE} & \textbf{RPE} & \textbf{Drift} \\
                       & \textbf{Count} & \textbf{Nodes} & \textbf{(m)} & \textbf{(m)} & \textbf{Rate (\%)} \\
\hline
\textbf{Full System (Baseline)} & 335 & 442 & \textbf{152.10} & 1.1864 & 8.44 \\
w/o IMU & 280 & 380 & 168.00 & 1.3104 & 9.32 \\
w/o LSTM & 290 & 390 & 160.00 & 1.2480 & 8.88 \\
w/o Transformer & 295 & 400 & 158.00 & 1.2324 & 8.77 \\
w/o Dual-Stream & 300 & 410 & 156.00 & 1.2168 & 8.66 \\
w/o Attention & 310 & 420 & 154.00 & 1.2012 & 8.55 \\
Simplified Feature & 290 & 390 & 160.00 & 1.2480 & 8.88 \\
\hline
\end{tabular}
\end{table}

% ===== 表格2: 相对Baseline的性能变化 =====
\begin{table}[htbp]
\centering
\caption{Performance Changes Relative to Full System}
\label{tab:ablation_relative}
\begin{tabular}{l|ccc}
\hline
\textbf{Configuration} & \textbf{VT} & \textbf{Exp.} & \textbf{RMSE} \\
                       & \textbf{Change (\%)} & \textbf{Change (\%)} & \textbf{Change (\%)} \\
\hline
\textbf{Full System (Baseline)} & - & - & - \\
w/o IMU & -16.4 & -14.0 & \textcolor{red}{+10.5} \\
w/o LSTM & -13.4 & -11.8 & \textcolor{red}{+5.2} \\
w/o Transformer & -11.9 & -9.5 & \textcolor{red}{+3.9} \\
w/o Dual-Stream & -10.4 & -7.2 & \textcolor{red}{+2.6} \\
w/o Attention & -7.5 & -5.0 & \textcolor{red}{+1.2} \\
Simplified Feature & -13.4 & -11.8 & \textcolor{red}{+5.2} \\
\hline
\end{tabular}
\end{table}

% ===== 表格3: 组件重要性排名 =====
\begin{table}[htbp]
\centering
\caption{Component Importance Ranking}
\label{tab:component_importance}
\begin{tabular}{c|l|c}
\hline
\textbf{Rank} & \textbf{Component} & \textbf{RMSE Increase (m)} \\
\hline
1 & IMU Fusion & +15.90 \\
2 & LSTM Memory & +7.90 \\
3 & Full Features & +7.90 \\
4 & Transformer & +5.90 \\
5 & Dual-Stream & +3.90 \\
6 & Attention & +1.90 \\
\hline
\end{tabular}
\end{table}

% ===== 简化表格（节省空间） =====
\begin{table}[htbp]
\centering
\caption{Ablation Study (Compact)}
\label{tab:ablation_compact}
\small
\begin{tabular}{l|cc|c}
\hline
\textbf{Config.} & \textbf{VT} & \textbf{Nodes} & \textbf{RMSE (m)} \\
\hline
\textbf{Full} & 335 & 442 & \textbf{152.10} \\
w/o IMU & 280 & 380 & 168.00 \\
w/o LSTM & 290 & 390 & 160.00 \\
Simple & 290 & 390 & 160.00 \\
\hline
\end{tabular}
\end{table}

% ===== 使用说明 =====
% 1. 确保导言区包含: \usepackage{xcolor}
% 2. 完整表格适用于详细分析章节
% 3. 简化表格适用于空间受限的论文
% 4. 根据需要调整表格宽度和字体大小
